\chapter{GPU-Accelerated Bayesian Inference for a Five-Species Biofilm}
\label{ch:heine}

This chapter develops the \emph{in-vitro} pillar of the thesis: the Bayesian
identification of the species interaction matrix of a defined five-species oral
biofilm from the controlled time-course experiments of Heine et al.\
\cite{Heine2025PeriImplant}. It is the methodological origin of the work: the
Hamilton $\bphi$--$\bpsi$ model and the TMCMC machinery introduced in
Chapter~\ref{ch:theory} are here fitted to real data and validated against
measurements that never entered the calibration. The last section carries the
same model to sparse in-vivo data with a metabolic sign prior
(Section~\ref{sec:heine_signprior}), and Chapter~\ref{ch:ansys} brings the
point model into a finite element framework.

\medskip\noindent\textit{This chapter is based on the manuscript:} Nishioka K.,
Klempt F., Geisler H., Soleimani M., Heine N., Doll-Nikutta K., Mukherjee R.,
Szafra\'{n}ski S.P., Stiesch M., Junker P., ``GPU-accelerated Bayesian inference
of multi-species biofilm interaction parameters via TMCMC'' (in preparation).

%==============================================================================
\section{Experimental System}
\label{sec:heine_data}
%==============================================================================

Heine et al.~\cite{Heine2025PeriImplant} cultivated a defined consortium of five
representative oral taxa, \textit{Streptococcus oralis} (\textit{So}),
\textit{Actinomyces naeslundii} (\textit{An}), \textit{Veillonella} spp.\
(\textit{Vei}), \textit{Fusobacterium nucleatum} (\textit{Fn}) and
\textit{Porphyromonas gingivalis} (\textit{Pg}), spanning the ecological
succession from early aerobic colonisers to the late anaerobic pathogen. The
experiment is organised as a $2\times2$ design crossing \emph{community state}
(commensal vs.\ dysbiotic) with \emph{cultivation mode} (static vs.\ HOBIC flow
reactor), giving the four conditions used throughout this thesis: Commensal
Static (CS), Commensal HOBIC (CH), Dysbiotic Static (DS) and Dysbiotic HOBIC
(DH).

The two community states employ different \textit{Pg} strains:
the avirulent \textit{Pg}~DSM\,20709 in the commensal consortium and the
virulent \textit{Pg}~W83, which carries a full complement of gingipain and fimbrial
virulence determinants, in the dysbiotic consortium. The dysbiotic/HOBIC
combination therefore reproduces, in a reductionist setting, the shear-driven
pathogen succession implicated in peri-implantitis.

Relative species abundances were quantified at six time points (Days 1, 3, 6, 10,
15, 21) with $N=8$ biological replicates per condition, normalised to unit sum at
each time point; Day~1 is consumed as the initial condition, leaving $N_t=5$
usable time points. Membrane-intact (viability) ratios $\psi_i^{\mathrm{exp}}$ and
HOBIC pH time series were measured independently and provide the held-out
validation targets of Section~\ref{sec:heine_validation}. With at most $25$
abundance observations against $15$ interaction unknowns, the inference operates
in the sparse-data regime that motivates the Bayesian, prior-regularised
treatment.

%==============================================================================
\section{Five-Species Inference Setup}
\label{sec:heine_setup}
%==============================================================================

The forward model is the coupled $\bphi$--$\bpsi$ Hamilton system of
Section~\ref{sec:theory_hamilton} integrated by implicit Euler with Newton's
method ($\Delta t=10^{-4}$, $2{,}500$ steps over the normalised window
$t_{\max}=0.25$ representing the 21-day experiment). The gauge $c^*=25$ fixes the
otherwise non-identifiable nutrient scale (Section~\ref{sec:theory_replicator});
since the assays use no antibiotic, $\alpha^*=0$ deactivates the decay term and
the sensitivity matrix $\mathbf{B}$. What remains to be inferred is the symmetric
interaction matrix $\mathbf{A}$.

\subsection{The 15-parameter interaction matrix}
\label{sec:heine_matrix}

The variational derivation forces $A_{ij}=A_{ji}$, reducing the $5\times5$ matrix
to $15$ free entries: $5$ diagonal self-interactions and $10$ symmetric
off-diagonals. I collect them into the parameter vector
\begin{equation}
  \btheta =
  \underbrace{(a_{11},a_{12},a_{22})}_{\text{\textit{So}--\textit{An}}}
  \oplus \underbrace{(a_{33},a_{34},a_{44})}_{\text{\textit{Vei}--\textit{Fn}}}
  \oplus \underbrace{(a_{13},a_{14},a_{23},a_{24})}_{\text{cross}}
  \oplus \underbrace{(a_{15},a_{25},a_{35},a_{45})}_{\text{\textit{Pg} cross}}
  \oplus \underbrace{(a_{55})}_{\text{\textit{Pg}}}
  \;\in\;\Real^{15},
  \label{eq:heine_blocks}
\end{equation}
whose five biologically motivated blocks structure the posterior analysis below.
Heine et al.\ report five experimentally confirmed interaction pathways
(Table~\ref{tab:heine_network}); for the remaining five pairs no direct
mechanism is documented. Rather than imposing those absences as hard constraints
($\theta_k\equiv0$), \emph{all} $15$ parameters are estimated and the posterior is
left to decide which interactions the data support, a full-discovery strategy
whose recovery of the expected commensal sparsity (Section~\ref{sec:heine_results})
is an emergent result rather than an assumption.

\begin{table}[htbp]
  \centering\small
  \begin{tabular}{@{}lll@{}}
    \toprule
    \textbf{Pair} & \textbf{Mechanism} & \textbf{Type}\\
    \midrule
    \textit{So}--\textit{An}  & Co-aggregation     & Bidirectional\\
    \textit{So}--\textit{Vei} & Lactate exchange   & Bidirectional\\
    \textit{So}--\textit{Fn}  & Formate symbiosis  & Bidirectional\\
    \textit{Vei}$\to$\textit{Pg} & pH elevation    & Unidirectional\\
    \textit{Fn}--\textit{Pg}  & Peptide supply     & Bidirectional\\
    \bottomrule
  \end{tabular}
  \caption[Experimentally confirmed species interactions]{Experimentally
  confirmed species interactions~\cite{Heine2025PeriImplant}. All $15$ parameters
  are estimated without a priori locking; the posterior concentrates near zero for
  the five biologically inactive pairs.}
  \label{tab:heine_network}
\end{table}

\subsection{Likelihood and observation noise}
\label{sec:heine_likelihood}

Assuming independent Gaussian errors across species and time, the likelihood is
\begin{equation}
  p(\yobs\mid\btheta)
  = \prod_{k=1}^{N_t}\prod_{i=1}^{5}
    \frac{1}{\sqrt{2\pi}\,\sigma_i}
    \exp\!\Bigl(-\tfrac{(y_{\mathrm{obs},i}(t_k)-\hat{y}_i(t_k;\btheta))^2}{2\sigma_i^2}\Bigr),
  \label{eq:heine_like}
\end{equation}
where $\sigma_i$ is a species-specific noise fixed from the replicate standard
deviation of the normalised fractions ($N=8$). It ranges from
$\sigma_i\!\approx\!0.05$ (\textit{Fn}, \textit{Pg} in commensal states, where
fractions are near zero) to $\sigma_i\!\approx\!0.20$ (\textit{So}, the most
variable species), with a condition-averaged mean of
$\bar{\sigma}\!\approx\!0.11$--$0.12$. The prior is a product of independent
uniforms (Section~\ref{sec:theory_tmcmc}) whose bounds are initialised broadly and
iteratively narrowed.

\subsection{Two-phase fix-$\psi$ strategy}
\label{sec:heine_twophase}

Direct joint estimation of $(\mathbf{A},\bpsi)$ is defeated by the multimodality
of the coupled $\bphi$--$\bpsi$ dynamics. I break the coupling in two phases.
In Phase~1 the viability $\psi_i(t_k)$ is fixed to the measured
membrane-intact ratios (fix-$\psi$), reducing the Jacobian dimension from $2n+2$ to
$n+2$ and removing the $\psi$-related modes; an initial wide-prior run furnishes
posterior quantiles and the production run uses the narrowed bounds
$[q_{0.01}-1.5\sigma,\,q_{0.99}+1.5\sigma]$. In Phase~2 the $\psi_i$ are
released as free variables, the viability data enter the likelihood as a soft
constraint, and the Phase-1 posterior serves as an informative prior on
$\mathbf{A}$, yielding a self-consistent joint posterior. Comparing the Phase-1
and Phase-2 MAP estimates (Section~\ref{sec:heine_discussion}) is itself a
diagnostic: agreement validates the fix-$\psi$ approximation, divergence localises
the $\psi$--$\mathbf{A}$ coupling.

\subsection{GPU-accelerated forward model}
\label{sec:heine_gpu}

Each production run requires $\mathcal{O}(10^5)$--$\mathcal{O}(10^6)$ forward ODE
solves ($N_p$ up to $10{,}000$ particles $\times$ $K=20$ mutations $\times$
$M\!\approx\!5$ TMCMC stages). The implicit-Euler Newton solver, including the Jacobian
assembly and the linear solve, is compiled with \texttt{jax.jit} into a
single fused XLA kernel and the entire particle ensemble is evaluated in parallel
via \texttt{jax.vmap} on an NVIDIA RTX\,4090. On a matched benchmark (DH,
$N_p=500$, $9$ stages) this delivers a $147\times$ end-to-end
speedup ($170\times$ per stage) over the 12-core CPU baseline (Table~\ref{tab:heine_timing}); at the
production scale of $N_p=10{,}000$ the per-condition wall-clock is
$1.7$--$2.6$\,h. The gradient-based samplers behind general probabilistic-programming
frameworks~\cite{Carpenter2017Stan, Phan2019NumPyro, Lao2020TFP} (HMC/NUTS) are ill-suited here:
variable-length leapfrog trajectories cause warp divergence that erases the
\texttt{vmap} advantage.

\begin{table}[htbp]
  \centering\small
  \begin{tabular}{@{}lrrr@{}}
    \toprule
    \textbf{Metric} & \textbf{CPU (12 cores)} & \textbf{GPU (RTX\,4090)} & \textbf{Speedup}\\
    \midrule
    \multicolumn{4}{@{}l}{\textit{Matched benchmark (DH, $N_p=500$, $9$ stages):}}\\
    \quad Per stage & $2{,}400$\,s & $14$\,s & $170\times$\\
    \quad Total     & $21{,}600$\,s & $147$\,s & $147\times$\\
    \midrule
    \multicolumn{4}{@{}l}{\textit{Production ($4$ cond., $N_p=10{,}000$, Phase~2 free $\psi$):}}\\
    \quad Per condition & $\sim$$120$\,h$^\dagger$ & $1.7$--$2.6$\,h & $\sim$$50\times$\\
    \quad All $4$ cond. & $\sim$$480$\,h$^\dagger$ & $8$\,h & $\sim$$60\times$\\
    \bottomrule
  \end{tabular}
  \caption[Computational cost comparison]{Computational cost. The matched
  $500$-particle benchmark isolates the GPU speedup at $147\times$; the
  production speedup is lower ($\sim$$50$--$60\times$) because $20\times$ more
  particles saturate GPU occupancy. $^\dagger$Extrapolated from the CPU benchmark
  assuming linear scaling in $N_p$.}
  \label{tab:heine_timing}
\end{table}

%==============================================================================
\section{Results}
\label{sec:heine_results}
%==============================================================================

The two-phase procedure is run for all four conditions with a CoV target
$\delta_{\mathrm{target}}=1.0$ for the adaptive tempering schedule $\tau_m$, proposal scale
$s^2=2.38^2/d$ (the optimal-scaling heuristic for random-walk
Metropolis~\cite{Roberts2001Optimal}), and $K=20$ mutation steps; $s$ is adapted online toward
the optimal acceptance rate of about $23\%$. The Phase-2
($N_p=10{,}000$) posteriors are the canonical product used throughout the thesis.

\subsection{Quantitative fit}
\label{sec:heine_fit}

Table~\ref{tab:heine_rmse} reports both phases. Phase~1 (fix-$\psi$, perfect
viability) is the upper bound on fit quality; Phase~2 is the physically consistent
model. The RMSE increase from Phase~1 to Phase~2 is at most $0.033$ (CH), confirming the
interaction structure is robust to releasing $\psi$. DS attains the best fit
(RMSE $0.033$) and DH retains strong pathogen identification
($R^2_{\textit{Pg}}=0.90$) despite the added viability degrees of freedom
(per-species $R^2$ in Table~\ref{tab:heine_r2}). In commensal states the
\textit{Fn}/\textit{Pg} fractions are near zero throughout, so their $R^2$ is
ill-defined ($\mathrm{SS}_{\mathrm{tot}}\!\approx\!0$, marked ``$-$''); the MAP
prediction there is $<\!0.01$, within the noise. \textit{So} shows negative $R^2$
in both phases except CH, which reflects its high replicate variability
($\sigma_i\!\approx\!0.20$).

\begin{table}[htbp]
  \centering\small
  \begin{tabular}{@{}l ccccc ccccc@{}}
    \toprule
    & \multicolumn{5}{c}{\textbf{Phase 1: fix-$\psi$ ($N_p\!=\!1{,}000$)}}
    & \multicolumn{5}{c}{\textbf{Phase 2: free $\psi$ ($N_p\!=\!10{,}000$)}}\\
    \cmidrule(lr){2-6}\cmidrule(lr){7-11}
    \textbf{Cond.} & RMSE & MAE & $a_{\mathrm{acc}}$ & $\max\ln\mathcal{L}$ & $\ln\hat Z$
                   & RMSE & MAE & $a_{\mathrm{acc}}$ & $\max\ln\mathcal{L}$ & $\ln\hat Z$\\
    \midrule
    CS & $0.119$ & $0.077$ & $0.14$ & $-6.3$ & $-2.6$ & $0.119$ & $0.075$ & $0.12$ & $-21.6$ & $-3.3$\\
    CH & $0.071$ & $0.053$ & $\mathbf{0.17}$ & $-3.6$ & $\mathbf{-1.3}$ & $0.104$ & $0.088$ & $0.07$ & $-112.3$ & $-8.5$\\
    DS & $\mathbf{0.020}$ & $\mathbf{0.025}$ & $0.11$ & $\mathbf{-1.2}$ & $-3.8$ & $\mathbf{0.033}$ & $\mathbf{0.024}$ & $\mathbf{0.13}$ & $-93.4$ & $\mathbf{-2.8}$\\
    DH & $0.067$ & $0.045$ & $0.10$ & $-6.6$ & $-9.6$ & $0.087$ & $0.060$ & $0.07$ & $-170.4$ & $-11.9$\\
    \bottomrule
  \end{tabular}
  \caption[Two-phase estimation results]{Two-phase estimation. $a_{\mathrm{acc}}$: mean
  MH acceptance rate. $\max\ln\mathcal{L}$ is not comparable across phases (Phase~2 adds
  a viability likelihood channel); use RMSE for cross-phase comparison.}
  \label{tab:heine_rmse}
\end{table}

\begin{table}[htbp]
  \centering\small
  \begin{tabular}{@{}l ccccc ccccc@{}}
    \toprule
    & \multicolumn{5}{c}{\textbf{Phase 1 (fix-$\psi$)}} & \multicolumn{5}{c}{\textbf{Phase 2 (free $\psi$)}}\\
    \cmidrule(lr){2-6}\cmidrule(lr){7-11}
    \textbf{Cond.} & \textit{So} & \textit{An} & \textit{Vei} & \textit{Fn} & \textit{Pg}
                   & \textit{So} & \textit{An} & \textit{Vei} & \textit{Fn} & \textit{Pg}\\
    \midrule
    CS & $-0.39$ & $0.80$ & $-3.2$ & $-$ & $-$ & $-0.36$ & $\mathbf{0.91}$ & $-3.7$ & $-$ & $-$\\
    CH & $0.55$ & $0.70$ & $0.59$ & $-$ & $-$ & $0.40$ & $-0.22$ & $-0.14$ & $-$ & $-$\\
    DS & $-0.87$ & $0.76$ & $\mathbf{0.93}$ & $0.54$ & $\mathbf{0.85}$ & $-0.86$ & $0.73$ & $\mathbf{0.91}$ & $0.66$ & $\mathbf{0.85}$\\
    DH & $-0.47$ & $\mathbf{0.83}$ & $\mathbf{0.83}$ & $0.60$ & $\mathbf{0.92}$ & $-0.21$ & $\mathbf{0.86}$ & $0.63$ & $0.06$ & $\mathbf{0.90}$\\
    \bottomrule
  \end{tabular}
  \caption[Per-species $R^2$]{Per-species $R^2$. ``$-$'': fraction $<0.02$ at all
  time points ($R^2$ ill-defined). Bold: $R^2>0.8$. DS/DH retain strong
  \textit{Pg} identification in both phases.}
  \label{tab:heine_r2}
\end{table}

\paragraph{Convergence.} Phase~2 converges in $M=4$--$7$ stages (ESS $=N_p=10{,}000$);
two-seed Phase-1 runs satisfy $\hat R<1.1$ for all parameters. Acceptance rates
of $7$--$13\%$ reflect bounded, non-Gaussian support; the high ESS confirms adequate mixing.

\subsection{Posterior predictive fits}
\label{sec:heine_fits}
% 4 Oct 2026: "NUTS" here changed to "TMCMC": the RMSE values (0.119/0.104/
% 0.033/0.087) and N_p = 10,000 are those of the TMCMC Phase-2 posterior, and
% Sec. heine_gpu states that NUTS is not used. Confirm against the figure source.

To assess fit quality independently of the symmetry constraint, I also fit an
asymmetric gLV model (free $A_{ij} \neq A_{ji}$) by L-BFGS-B point estimation.
Table~\ref{tab:glv_hamilton_comparison} summarises the two models across four
axes; Figure~\ref{fig:heine_glv_fit} shows the gLV MAP fits.
The gLV achieves markedly lower RMSE (CS: 0.032, CH: 0.026, DS: 0.012,
DH: 0.022) compared to the Hamilton TMCMC posterior (CS: 0.119, CH: 0.104,
DS: 0.033, DH: 0.087), confirming that the symmetry constraint
$A_{ij}=A_{ji}$ limits trajectory accuracy.
However, gLV's apparent advantage partly reflects over-fitting:
with $N_t=5$ usable time points ($25$ observations) and $30$ free parameters, the
model has more free parameters than observations. Its out-of-sample error was not
assessed; for the Hamilton model, the cross-condition transfer RMSE
(CS$\to$CH $=0.098$ vs.\ CH self-prediction $=0.104$; Table~\ref{tab:heine_transfer})
stays close to the in-sample value.
The Hamilton posterior provides three capabilities that gLV cannot: uncertainty
quantification (credible bands, Figure~\ref{fig:heine_fits}), exact simplex
conservation, and the falsifiable pH prediction of Section~\ref{sec:heine_validation}.

\begin{table}[htbp]
  \centering\small
  \renewcommand{\arraystretch}{1.3}
  \begin{tabular}{@{}lcc@{}}
    \toprule
    & \textbf{gLV (L-BFGS-B)} & \textbf{Hamilton TMCMC} \\
    \midrule
    Interaction matrix $\bA$ & Asymmetric ($25$ entries) & Symmetric ($15$ entries) \\
    Parameters total         & 30 ($A$ + $b$)              & 20 ($A$ + $\psi$) \\
    RMSE range (all cond.)   & $0.012$--$0.032$            & $0.033$--$0.119$ \\
    Uncertainty              & Point estimate only         & Full posterior ($N_p=10{,}000$) \\
    Simplex $\sum_i\phi_i=1$ & Enforced post-hoc           & Exact (variational) \\
    Thermodynamic consistency & No                          & Yes (Lyapunov function) \\
    pH prediction (in vitro) & Not possible                & $R^2=0.778$ (held-out) \\
    \bottomrule
  \end{tabular}
  \caption[gLV vs.\ Hamilton model comparison]{Comparison of the asymmetric gLV
  (L-BFGS-B point estimate) and the Hamilton replicator (TMCMC posterior).
  Lower RMSE for gLV reflects higher model flexibility \emph{and} over-fitting
  ($25$ observations, 30 free parameters); the gLV out-of-sample error was not
  assessed (Table~\ref{tab:heine_transfer} covers the Hamilton model only).
  Hamilton's advantages are uncertainty quantification, exact compositional
  consistency, and probabilistic predictions of held-out chemistry.}
  \label{tab:glv_hamilton_comparison}
\end{table}

\begin{figure}[p]
  \centering
  \includegraphics[width=0.80\textwidth]{heine_glv_fit_thesis.pdf}
  \caption[gLV MAP fits across four conditions]{gLV MAP fits of volume fraction
  $\phi_i$. Layout identical to Figure~\ref{fig:heine_fits}: columns CS, CH,
  DS, DH; rows: five species. Solid line: MAP trajectory; circles: experimental
  median; boxes: inter-replicate IQR. \textit{V.~dispar} (commensal) vs.\
  \textit{V.~parvula} (dysbiotic) and \textit{Pg} strain (DSM\,20709 vs.\
  W83) are labelled within each panel. RMSE values are shown in the column
  titles.}
  \label{fig:heine_glv_fit}
\end{figure}

Figure~\ref{fig:heine_fits} shows the Hamilton TMCMC posterior predictive
trajectories. The model reproduces commensal homeostasis (CS:
\textit{So}/\textit{An} growth with \textit{Pg} near zero although it is not
fixed; CH: the \textit{So}-dominated bloom under flow), the
clean pathogen succession of DS, and the
non-linear \textit{Vei}/\textit{Pg} amplification of DH, the late-onset surge
driven by bridge-mediated cross-feeding. Credible bands widen from CS to DH
with the increasing effective dimensionality.

\begin{figure}[p]
  \centering
  \includegraphics[width=0.82\textwidth]{heine_fits_4cond.pdf}
  \caption[Hamilton TMCMC posterior predictive fits across four conditions]{Hamilton
  TMCMC posterior predictive fits of volume fraction $\phi_i$
  ($N_p=10{,}000$). Columns: CS, CH, DS, DH; rows: five species. Solid: MAP
  $\phi_i$; dashed: living fraction $\bar\phi_i=\phi_i\psi_i$. Bands:
  $50\%$ (dark) / $80\%$ (light) credible intervals from $100$ posterior
  samples. Circles: experimental median ($N=8$); boxes: inter-replicate IQR.
  \textit{Pg} uses different strains by state: avirulent DSM\,20709
  (commensal) vs.\ virulent W83 (dysbiotic).}
  \label{fig:heine_fits}
\end{figure}

\subsection{Inferred interaction matrices and posteriors}
\label{sec:heine_matrices}

The MAP matrices (Figure~\ref{fig:heine_heatmap}) reveal a progressive
reorganisation. In CS/CH the dominant entries are the \textit{So}--\textit{An} and
\textit{So}--\textit{Vei} blocks (the mutualistic early-coloniser core), with the
\textit{Fn} and \textit{Pg} rows near zero. The dysbiotic transition activates
strong positive \textit{Fn}--\textit{Pg} and self-interaction entries:
$\hat a_{55}^{\mathrm{DH}}(\textit{Pg}\,\text{self})=3.43$,
$\hat a_{45}^{\mathrm{DH}}(\textit{Fn}\!\to\!\textit{Pg})=5.63$,
$\hat a_{44}^{\mathrm{DH}}(\textit{Fn}\,\text{self})=4.45$, versus much weaker DS
values ($0.30$, $2.20$, $1.01$) and near-zero CS/CH. The posterior violins
(Figure~\ref{fig:heine_violin}) make the sparsity explicit: the
\textit{Pg}-related entries ($a_{15},a_{25},a_{35},a_{45},a_{55}$) concentrate
sharply at zero in CS/CH without being fixed there, and shift to broad positive
distributions in DS/DH. DS exhibits the narrowest posteriors overall, consistent
with its lowest RMSE.

\begin{figure}[htbp]
  \centering
  \includegraphics[width=0.80\textwidth]{heine_heatmap_A_4cond.pdf}
  \caption[MAP interaction matrices]{MAP interaction matrices $\mathbf{A}$
  (Phase~2). Red: cooperative ($A_{ij}>0$); blue: competitive ($A_{ij}<0$). In
  CS/CH the \textit{Fn}/\textit{Pg} rows and columns are near zero; these entries
  are estimated, not fixed. The dysbiotic transition activates cooperative
  cross-feeding in the \textit{Fn}--\textit{Pg} block.}
  \label{fig:heine_heatmap}
\end{figure}

The gLV model (no symmetry constraint) yields the same qualitative pattern
with substantially lower RMSE (Table~\ref{tab:glv_hamilton_comparison}).
Figure~\ref{fig:heine_glv_heatmap} shows the corresponding MAP interaction
matrices with annotated entry values. The condition-specific reorganisation is
preserved, while the asymmetry of $\bA$ reveals directed interactions
suppressed by the Hamilton symmetry constraint.

\begin{figure}[htbp]
  \centering
  \includegraphics[width=0.80\textwidth]{heine_glv_Amatrix_thesis.pdf}
  \caption[gLV MAP interaction matrices across four conditions]{gLV MAP
  interaction matrices $\mathbf{A}$ (asymmetric). Red: cooperative
  ($A_{ij}>0$); blue: competitive ($A_{ij}<0$). Annotated values indicate
  entries with $|A_{ij}|>0.04$. The CS/CH commensal core
  (\textit{So}--\textit{An}/\textit{Vei}; \textit{Vd} in the figure) and the DS/DH dysbiotic
  \textit{Fn}--\textit{Pg} block are consistent with the Hamilton result
  (Figure~\ref{fig:heine_heatmap}), but the asymmetric entries expose
  directionality absent from the symmetric prior.}
  \label{fig:heine_glv_heatmap}
\end{figure}

Two features emerge from the annotated entries.
% 4 Oct 2026: values re-read from figures/heine_glv_Amatrix_thesis.pdf; the
% earlier text paired a CS and a CH entry and called both CS/CH.
\emph{Magnitude contraction in dysbiosis}: the largest CS/CH entries reach
$|A_{ij}|\!\approx\!1.8$ (CS: $A_{\text{So,Vei}}\!=\!+1.83$, $A_{\text{Vei,An}}\!=\!+1.67$;
CH: $A_{\text{An,Vei}}\!=\!+1.76$), whereas no DS/DH entry exceeds $0.58$, consistent with a tightly coordinated
commensal network versus a weaker but more distributed dysbiotic attractor.
\emph{Directed asymmetry and late-coloniser activation}: in CS, $A_{\text{Vei,An}}\!=\!+1.67$
but $A_{\text{An,Vei}}\!=\!-0.02$, a one-directional benefit invisible to the symmetric model;
in DS/DH the \textit{Pg}-related entries activate (DS: $A_{\text{An,Pg}}\!=\!+0.49$; DH:
$A_{\text{Fn,Pg}}\!=\!+0.07$), consistent with the sign-enrichment analysis of
Section~\ref{sec:heine_signprior}.

\begin{figure}[htbp]
  \centering
  \includegraphics[width=0.82\textwidth]{heine_posterior_violin.pdf}
  \caption[Posterior distributions of the 15 interaction entries]{Posterior
  distributions of the $15$ interaction entries (Phase~2). Stars: MAP. The
  concentration of \textit{Pg}-related entries at zero in CS/CH confirms that
  biological sparsity emerges from the data without a priori locking.}
  \label{fig:heine_violin}
\end{figure}

\subsection{Structural reorganisation across conditions}
\label{sec:heine_reorg}

To quantify the inter-condition differences I compute, for each pair, the
correlation $\rho$ of the vectorised matrices and the relative Frobenius distance
$\Delta_F$ (Table~\ref{tab:heine_pairwise}). Matrices within the same health state
share a backbone ($\rho_{\text{CH--CS}}=+0.26$, $\rho_{\text{DH--DS}}=+0.40$),
whereas three of the four cross-state pairs anticorrelate ($\rho_{\text{DH--CS}}=-0.49$),
and all cross-state pairs have large distances ($\Delta_F\geq1.01$). A UMAP
embedding~\cite{McInnes2018UMAP} (cf.\ t-SNE~\cite{VanDerMaaten2008tSNE}) of the full posterior
(Figure~\ref{fig:heine_umap}) confirms four well-separated clusters with no
overlap. The transition from health to disease thus changes the sign pattern of
the interaction network (Table~\ref{tab:heine_rewiring}).

\begin{table}[htbp]
  \centering\small
  \begin{tabular}{@{}cc|ccc@{}}
    \toprule
    \textbf{C1} & \textbf{C2} & $\rho$ & $\Delta_F$ & $d_{15\mathrm{D}}$\\
    \midrule
    CH & CS & $+0.26$ & $1.06$ & $\mathbf{3.38}$\\
    DH & DS & $\mathbf{+0.40}$ & $\mathbf{0.88}$ & $5.86$\\
    CH & DH & $+0.24$ & $1.01$ & $6.33$\\
    CH & DS & $-0.27$ & $1.27$ & $4.88$\\
    DH & CS & $\mathbf{-0.49}$ & $1.28$ & $\mathbf{7.77}$\\
    CS & DS & $-0.27$ & $\mathbf{1.36}$ & $5.88$\\
    \bottomrule
  \end{tabular}
  \caption[Pairwise comparison of MAP matrices]{Pairwise comparison of MAP
  matrices. $\rho$: correlation of $\mathrm{vec}(\mathbf{A})$; $\Delta_F$: relative
  Frobenius distance; $d_{15\mathrm{D}}$: posterior-centroid distance in 15D.}
  \label{tab:heine_pairwise}
\end{table}

\begin{figure}[htbp]
  \centering
  \includegraphics[width=0.56\textwidth]{heine_umap_A_3d.pdf}
  \caption[UMAP embedding of posterior A-matrix samples]{UMAP 3D embedding of the
  posterior A-matrix samples ($N_p=10{,}000$). The four conditions form
  well-separated clusters; centroid distances mirror the $d_{15\mathrm{D}}$ ordering
  (CS--CH closest, CS--DH most distant).}
  \label{fig:heine_umap}
\end{figure}

Figure~\ref{fig:heine_umap_comparison} overlays the gLV MAP solutions on the same
UMAP embedding (fitted on the Hamilton posterior; gLV projected via
$\tilde{A}_{ij}=(A_{ij}+A_{ji})/2$). Hamilton MAP stars ($\bigstar$) fall within their
posterior clouds in all four conditions; gLV MAP crosses ($\times$) deviate substantially
for CS and DH, migrating toward the CH region and the embedding centre respectively.
The lower gLV RMSE shows that the displacement reflects asymmetry unlocking regions of interaction space geometrically
inaccessible under the $A_{ij}=A_{ji}$ constraint.

\begin{figure}[htbp]
  \centering
  \includegraphics[width=0.64\textwidth]{heine_umap_comparison.pdf}
  \caption[UMAP comparison: Hamilton posterior and gLV MAP]{UMAP of the
  Hamilton posterior A-matrix samples (a subsample of $1{,}200$ per condition from
  the $N_p=10{,}000$ posterior; coloured clouds). Stars ($\bigstar$): Hamilton MAP; crosses ($\mathbf{\times}$):
  gLV MAP, projected via $\tilde{A}_{ij}=(A_{ij}+A_{ji})/2$. Dashed lines
  connect Hamilton MAP to gLV MAP of the same condition. Legend at right.
  CS and DH gLV MAP solutions are displaced well outside their Hamilton
  posterior clouds, indicating that asymmetry unlocks qualitatively different
  regions of interaction space.}
  \label{fig:heine_umap_comparison}
\end{figure}

Cross-condition \emph{transfer} sharpens the point (Table~\ref{tab:heine_transfer}):
evaluating each MAP on every other condition's data, the within-state commensal
transfer CS$\to$CH (RMSE $0.098$) is better than the CH self-prediction ($0.104$),
which points to a shared commensal backbone across cultivation modes, whereas
cross-state transfers (DS$\to$CS $0.384$, CS$\to$DH $0.198$) degrade sharply.

\begin{table}[htbp]
  \centering\small
  \begin{tabular}{@{}lcccc@{}}
    \toprule
    Source $\backslash$ Target & CS & CH & DS & DH \\
    \midrule
    CS & \textbf{0.119} & 0.098 & 0.261 & 0.198 \\
    CH & 0.257 & \textbf{0.104} & 0.100 & 0.159 \\
    DS & 0.384 & 0.259 & \textbf{0.033} & 0.255 \\
    DH & 0.342 & 0.212 & 0.289 & \textbf{0.087} \\
    \bottomrule
  \end{tabular}
  \caption[Cross-condition prediction RMSE]{Cross-condition prediction RMSE.
  Diagonal (bold): self-prediction. CS$\to$CH ($0.098$) outperforms CH
  self-prediction ($0.104$), confirming a shared commensal backbone.}
  \label{tab:heine_transfer}
\end{table}

\paragraph{Effective dimensionality.} With $r_i=\sigma_{\mathrm{post},i}/\Delta_{\mathrm{prior}}$
($\Delta_{\mathrm{prior}}=6$), all $15$ parameters satisfy $r_i<0.43$ in every
condition (DS tightest, mean $r=0.07$), and even the commensal \textit{Pg} entries
reach $r_i\leq0.10$: the data constrain them to near zero. The full $15$-parameter model is therefore not over-parametrised
(the identifiability diagnostic of Section~\ref{sec:theory_tmcmc}).

\subsection{Edge rewiring between health states}
\label{sec:heine_rewiring}

The structural reorganisation of Section~\ref{sec:heine_reorg} can be localised to
individual edges by classifying each pair as facilitative or competitive
from the full posterior (the facilitation probability $P_{\mathrm{fac}}$ is the
posterior mass with $A_{ij}>0$). A handful of edges flip sign with near-certainty
between the commensal and dysbiotic states (Table~\ref{tab:heine_rewiring}). The
\textit{So}--\textit{Vei} lactate axis, the metabolic core
of commensal health, reverses from strong mutual facilitation
($P_{\mathrm{fac}}\!\approx\!1.0$) to competition ($P_{\mathrm{fac}}\!<\!0.05$),
while \textit{F.~nucleatum} acquires facilitation from \textit{An} and
\textit{Vei} ($P_{\mathrm{fac}}$ rising from $0.09$--$0.23$ to $0.77$--$0.97$).
Consistently, the net outgoing effect of \textit{P.~gingivalis} is most negative in
dysbiotic-HOBIC ($\overline{\Delta}_{\textit{Pg}}^{\,\mathrm{DH}}=-0.14$,
$P_{\mathrm{fac}}=0.27$): in the surge regime \textit{Pg} is a net competitor
that consumes the shared niche rather than feeding it.

\begin{table}[htbp]
  \centering\small
  \begin{tabular}{@{}llcc@{}}
    \toprule
    Edge & Commensal $\to$ Dysbiotic & $P_{\mathrm{fac}}^{\mathrm{comm}}$ & $P_{\mathrm{fac}}^{\mathrm{dys}}$\\
    \midrule
    \textit{So}--\textit{Vei} & facilitation $\to$ competition & $1.00$ & $0.00$--$0.03$\\
    \textit{An}--\textit{Fn}  & competition $\to$ facilitation & $0.23$ & $0.97$\\
    \textit{Vei}--\textit{Fn} & competition $\to$ facilitation & $0.09$ & $0.77$\\
    \bottomrule
  \end{tabular}
  \caption[Posterior edge rewiring between health states]{Edges that flip sign with
  near-certainty between commensal and dysbiotic states (posterior facilitation
  probability $P_{\mathrm{fac}}$, from the $10{,}000$-particle posterior). $\mathbf A$
  is symmetric, so each edge is undirected; where two values are given they span the
  dysbiotic estimates. The
  commensal \textit{So}--\textit{Vei} lactate mutualism reverses to competition; the
  bridging organism \textit{F.~nucleatum} gains facilitation.}
  \label{tab:heine_rewiring}
\end{table}

%==============================================================================
\section{Independent Validation}
\label{sec:heine_validation}
%==============================================================================

The Phase-2 MAP parameters are validated against measurements that never entered
the calibration.

\paragraph{pH prediction (LOO-$R^2=0.92$; independent $R^2=0.78$).}
A five-predictor linear pH model, selected by leave-one-out CV on $N=90$ HOBIC
time points,
\begin{equation}
  \mathrm{pH}(t) = 6.10 + 0.16\,\phi_{\textit{So}} + 0.55\,\phi_{\textit{An}}
    + 0.30\,\phi_{\textit{Vei}} + 1.22\,\phi_{\textit{Fn}} + 1.88\,\phi_{\textit{Pg}},
  \label{eq:heine_ph}
\end{equation}
has its coefficients fixed from \emph{experimental} fractions alone
(LOO-$R^2=0.920$). To test the \emph{inferred} $\mathbf{A}$, $500$ Phase-2
posterior samples are integrated forward from Day-1 fractions and the fixed
coefficients applied to the ODE-predicted fractions; no pH data enter
the TMCMC. This gives an independent $R^2=0.778$ (RMSE $0.126$\,pH units)
against the held-out HOBIC pH series (Figure~\ref{fig:heine_ph}), confirming that
$\mathbf{A}$ captures the community drivers of pH. The species--pH correlations are
mechanistically coherent: \textit{Fn} ($r=+0.93$) and \textit{Pg} ($r=+0.93$)
dominate alkalinisation (gingipain-mediated arginine/lysine catabolism releasing
ammonia), while \textit{So} acidifies ($r=-0.68$) via lactate.

\begin{figure}[htbp]
  \centering
  \includegraphics[width=0.80\textwidth]{heine_ph_validation.pdf}
  \caption[Independent pH validation]{Independent pH validation. \textbf{(a)}
  Measured HOBIC pH (solid) vs.\ TMCMC posterior predictions (dashed, markers) with
  $90\%$ credible bands, commensal (blue) and dysbiotic (red). \textbf{(b)}
  Predicted vs.\ measured pH at six time points. Regression coefficients were
  fitted from species fractions only (LOO-CV, $N=90$); pH was not used in
  calibration. $R^2=0.778$, RMSE $0.126$\,pH units.}
  \label{fig:heine_ph}
\end{figure}

\paragraph{Further checks.} The predicted DH \textit{Pg} trajectory correlates
with the measured gingipain profile ($r=0.90$), both showing delayed Day-15--21
onset. The signs of the posterior are metabolically plausible: the \textit{So}--\textit{Vei}
lactate axis is facilitative in the commensal states and competitive in the
dysbiotic states, and the \textit{Fn} edges gain facilitation in the dysbiotic
states (Table~\ref{tab:heine_rewiring}).
Section~\ref{sec:heine_signprior} uses such metabolic signs as a prior for
in-vivo data.

%==============================================================================
\section{Discussion}
\label{sec:heine_discussion}
%==============================================================================

\paragraph{Biological interpretation.} The inferred matrices give a quantitative
account of the commensal-to-dysbiotic transition. The healthy state rests on a
\textit{So}--\textit{An} co-aggregation and \textit{So}--\textit{Vei} lactate core;
dysbiosis activates cooperative \textit{Vei}$\to$\textit{Pg} and
\textit{Fn}$\to$\textit{Pg} pathways. That structural similarity is preserved
within health states ($\rho_{\text{within}}>0$) but lost across them
($\rho_{\text{cross}}<0$ for three of the four cross-state pairs) shows the transition to be a network reorganisation
rather than a magnitude shift, and the full-discovery posterior recovers commensal
sparsity as an emergent property.

\paragraph{A falsifiable prediction.} Because the terminal \textit{Pg} surge is
gated by the cooperative \textit{Fn}$\to$\textit{Pg} coefficient, which is
indistinguishable from zero in both commensal states but rises to
$\hat a_{45}^{\mathrm{DH}}=5.63$ (posterior mean $+4.86$, $95\%$ CI
$[+3.03,+5.95]$) under dysbiotic flow, the model predicts that
\emph{attenuating this single interaction}, e.g.\ by omitting
\textit{F.~nucleatum} from the consortium, should suppress the terminal \textit{Pg}
bloom and drive the DH trajectory toward a commensal-like, \textit{Pg}-low steady
state. This is a direct leave-one-species-out co-culture test that would
discriminate the network hypothesis from a \textit{Pg}-intrinsic growth
explanation, and it recapitulates the established bridging role of
\textit{F.~nucleatum}~\cite{Kolenbrander2006FnBridge, Signat2011Fn}.

\paragraph{Fix-$\psi$ approximation and Phase~2.} Figure~\ref{fig:heine_phase}
compares the Phase-1 and Phase-2 MAPs parameter by parameter. In commensal
conditions they agree closely ($r=0.83$--$0.90$): the $\psi$--$\mathbf{A}$ coupling
is weak and fix-$\psi$ is a good approximation. In dysbiotic conditions they
diverge ($r=0.23$ for DS, $0.62$ for DH), with the largest shifts in
\textit{Pg}-related entries, where viability dynamics are most active.
This condition-dependent consistency validates the two-phase design: Phase~1 is a
reliable starting point, and Phase~2 corrections matter where $\psi$ couples
non-negligibly to $\mathbf{A}$. The viability predictions themselves remain poor
(RMSE $0.11$--$0.39$); the $\psi$ channel functions as a soft constraint that
breaks the $\psi$--$\mathbf{A}$ degeneracy rather than as a predictive model of
membrane integrity.

\begin{figure}[htbp]
  \centering
  \includegraphics[width=0.80\textwidth]{heine_phase1_vs_phase2.pdf}
  \caption[Phase~1 vs.\ Phase~2 MAP estimates]{Phase~1 (fix-$\psi$) vs.\ Phase~2
  (free $\psi$) MAP estimates; each point is one of the $15$ parameters, dashed
  line is identity. Commensal conditions agree strongly ($r\geq0.83$); dysbiotic
  conditions shift systematically in \textit{Pg}-related parameters, confirming the
  necessity of Phase~2.}
  \label{fig:heine_phase}
\end{figure}

\paragraph{Model evidence and computation.} The log-evidence $\ln\hat Z$~\cite{Kass1995Bayes}
(Table~\ref{tab:heine_rmse}) is higher for the commensal conditions in Phase~1 under the same
$15$-parameter prior, because the posterior concentrates pathogen entries near
zero and so reduces the Occam factor, a model-based signature of ecological
complexity. Methodologically, the $147\times$ GPU speedup makes
$10{,}000$-particle posteriors routine and is what renders the full-discovery,
no-locking strategy tractable in $15$ dimensions. These in-vitro interaction
matrices, and the validation that they predict held-out chemistry, calibrate the
point model that Chapter~\ref{ch:ansys} brings into space.

\paragraph{Limitations and the symmetry approximation.}
Three limitations bound interpretation of the in-vitro matrices. First, the
Hamilton symmetry constraint $A_{ij}=A_{ji}$ excludes directed
contact-dependent adhesion and toxin delivery by construction. Second, the
five-species consortium is a curated subset; unmeasured community members
could absorb metabolites that shift apparent cross-feeding signs. Third,
the Phase-1 posteriors are conditioned on the measured viability, and Phase~2
treats $\psi$ only as a soft constraint; a fully
joint $(\mathbf{A},\psi)$ inference would propagate additional uncertainty
into the interaction estimates, particularly in dysbiotic conditions where
the Phase-2 correction is largest. These approximations are consistent with
the in-vitro setting and motivate the in-vivo cross-validation of
Section~\ref{sec:heine_signprior}, where the symmetry constraint is
re-examined against a ten-patient longitudinal dataset.

%==============================================================================
\section{From In-Vitro to In-Vivo Data: a Metabolic Sign Prior}
\label{sec:heine_signprior}
%==============================================================================
% Condensed 3 Oct 2026 from the former chapter 4 (chapters/ch4_dieckow.tex,
% kept in the tree for its full text and figures). Internship work.

The calibration above rests on rich in-vitro data: five species, six time points,
four controlled conditions. Clinical data are much sparser. This section shows,
in condensed form, how the same type of model can still be identified there when
the signs of the interactions are fixed beforehand from metabolic models. The
work was done during an internship at NIFE together with S.\,P.\ Szafra\'{n}ski
and is summarised here because it completes the calibration side of the
thesis; the full analysis is available as an internal report.

\paragraph{Data.} Dieckow et al.~\cite{dieckow2024structure} characterised
early biofilm formation on implant abutments in $12$ patients. The analysis
uses $10$ of them, each sampled at three weekly time points. The
16S~rRNA amplicon reads were aggregated into $10$ class-level guilds and
normalised to the simplex, $\sum_i\phi_i=1$. With $10$ guilds and three time
points, the first of which is the initial condition, there are $20$ observations per patient against up to $55$ entries of a
symmetric interaction matrix, so the inverse problem cannot be solved without
additional information.

\paragraph{Sign prior.} The prior of Section~\ref{sec:theory_fba} is assembled
in three layers: experimentally documented production and use of metabolites
from eHOMD~\cite{Escapa2018eHOMD} and Szafra\'{n}ski et
al.~\cite{Szafranski2025} ($34$ constrained guild pairs); parsimonious FBA of
one AGORA2 model per guild~\cite{Heinken2023AGORA2} ($35$ pairs); and
community FBA with MICOM~\cite{Diener2020MICOM} ($72$ pairs). The directed
signals are symmetrised for the symmetric model, and the penalty acts only on
coefficients of the wrong sign.

\paragraph{Model and test.} The replicator form of
Section~\ref{sec:theory_replicator} is fitted with one interaction matrix
$\bA$ shared by all patients and a growth vector $\bb^{(p)}$ per patient. It is
tested by leave-one-patient-out cross-validation: $\bA$ is fitted on nine
patients and frozen, the held-out patient's $\bb^{(p)}$ is fitted from the
first week only, and the second and third weeks are predicted.
Table~\ref{tab:heine_signprior} lists the prediction errors.

\begin{table}[htbp]
  \centering\small
  \begin{tabular}{@{}llcc@{}}
    \toprule
    model & prior & sign agreement & LOO-RMSE\\
    \midrule
    replicator, symmetric $\bA$ & none            & --        & $0.0595$\\
    replicator, symmetric $\bA$ & eHOMD + pFBA    & $94\%$    & $0.0516$\\
    replicator, symmetric $\bA$ & + MICOM         & $100\%$   & $0.0502$\\
    gLV, asymmetric             & none            & --        & $0.0588$\\
    gLV, asymmetric             & eHOMD + pFBA    & $97\%$    & $0.0490$\\
    Bayesian gLV (MDSINE2~\cite{Gibson2021MDSINE2}) & sparsity & -- & $0.0552$\\
    \bottomrule
  \end{tabular}
  \caption[Leave-one-patient-out prediction with a sign prior]{Leave-one-patient-out
    prediction of the in-vivo guild compositions ($10$ patients of Dieckow et al.).
    Sign agreement: share of the constrained pairs whose fitted sign matches the
    prior. The persistence null model has a Bray--Curtis dissimilarity of
    $0.281$; the best model reaches $0.147$.}
  \label{tab:heine_signprior}
\end{table}

\paragraph{Results.}
\begin{itemize}
  \item The sign prior improves the symmetric model step by step, from
        $0.0595$ without prior to $0.0502$ with all three layers. The
        asymmetric gLV with the first two layers is slightly better
        ($0.0490$); its directed coefficients add freedom that the symmetric
        model does not have. The Bayesian gLV collapses to $\bA\approx0$ with
        three time points.
  \item Only the sign carries over from the metabolic models. A prior that
        ties $A_{ij}$ to the magnitude of the predicted fluxes recovers the
        correct sign for only $4$ to $11$ of $70$ pairs. Fitted without any
        prior, the symmetric model still recovers $11$ of the $14$
        cross-feeding signs predicted by AGORA2 ($78.6\,\%$, against $37.7\,\%$
        under a permutation null, $p=4\times10^{-4}$); the asymmetric gLV
        shows no such agreement ($p=0.96$).
  \item On $127$ independent peri-implant samples of Joshi et
        al.~\cite{joshi2025}, a dysbiosis ratio of the guilds ordered by the
        model separates peri-implantitis from health and mucositis
        (Kruskal--Wallis $p<10^{-4}$; Spearman $\rho=0.46$ with severity).
\end{itemize}

\paragraph{Limits.} With $10$ patients the gain of the community-FBA layer over
the first two is only marginally significant ($p\approx0.07$). The symmetric
matrix cannot represent directed effects, which is where the asymmetric gLV
gains its edge. For the rest of this thesis the result that matters is a
narrower one: the signs of the interaction matrix are the part of the model
that is supported by independent evidence, in vitro and in vivo.
